\section{Introduction}
Engineers use logs, traces, and operational records to investigate outages.
LLM-based RCA systems also use this evidence
\cite{ahmed2023recommending,chen2024automatic,wang2024rcagent,xu2025openrca}.
The records can contain customer IDs, account links, request history, time
data, and secrets. Sending all records to an external model can expose data
that the task does not need. A generic redaction marker can also remove the
joins that explain the failure. A stable pseudonym keeps joins, but it can
expose frequency and linkage. Thus, the important question is which relations
remain visible after a transformation.

DiagPriv names these task-required relations \emph{diagnostic invariants}.
Examples include repeated request identity, identifier shape, and event order.
The study measures both RCA utility and information exposure. It does not
claim that transformed data is anonymous.

Three deterministic synthetic families test different dependencies.
Incident 001 needs identifier linkage. Incident 002 needs identifier shape.
Incident 003 v3 needs temporal order and overlap. Table~\ref{tab:families}
summarizes the designs.
\begin{table}[H]
\centering
\small
\resizebox{\textwidth}{!}{%
\begin{tabular}{llll}
\toprule
Family & Observed problem & Key invariant & Mechanism contrast \\
\midrule
001 & Duplicate payment effect & Identifier equality and cross-source linkage & Timeout, re-arrival, missing idempotency \\
002 & Missing provisioning effect & External-reference format and shape & Malformed reference, exact lookup miss \\
003 v3 & Premature compensation / twin & Deadline order and interval overlap & Overlapping execution versus serialized invalid transition \\
\bottomrule
\end{tabular}%
}
\caption{Three synthetic incident families and their intended diagnostic dependencies.}
\label{tab:families}
\end{table}


Prior work covers task-aware anonymization and evidence-grounded RCA
\cite{loiseau2025taueval,loiseau2026adaptive,wang2026cloudopsbench}. This work
adds a bounded evaluation design. It separates the oracle from the ModelView,
compares fixed transformations on matched development cases, reports an
exposure vector, and scores cited evidence separately from RCA labels. It does
not implement an automatic transformation selector.
